% Unregistered structural candidate. Full independent/root audits required.
\subsection{A single horizontal branch section in mixed characteristic}

\paragraph{Candidate local statement.}
Let \(R\) be a complete discrete valuation ring of mixed
characteristic \((0,p)\), with algebraically closed residue field,
and let \(X/R\) be a smooth proper geometrically connected curve
of genus at least one. Let \(x\in X_s\) be a closed point.
Suppose an actual finite Galois covering of \(X_K\), with
\(p\)-group \(G\), has at most one horizontal branch section
specializing to \(x\), and that this section is a smooth
\(R\)-section. After finite extension all the other branch
sections are assumed to have reductions different from \(x\).
Then every positive-genus component of a semistable source
contracted to \(x\) has \(p\)-rank zero.

This statement concerns contracted components. It does not
identify the genus of a marked vertical fibre or a joining
path between its marks. No slope assertion for an arbitrary
marked fibre is imported from Raynaud.

\begin{proof}[Candidate proof of the local statement]
If there is no branch section through \(x\), the assertion is
Raynaud's nonmarked smooth-point theorem, in the precise form
of Yang, Definition~1.5, Proposition~1.6 and Theorem~1.9
\cite{fghYang2024}. Suppose there is a branch section \(P\).
Make a finite extension to define its geometric inertia and
the roots of unity used below. The inertia at \(P_K\) is
cyclic, because this is a discrete valuation of residue
characteristic zero. Write its order as \(e=p^a\).

\paragraph{An actual global Kummer base change with a smooth local model.}
Choose a horizontal relatively ample divisor \(D\) disjoint
from \(P\) and from \(x\). For a sufficiently large multiple
of \(D\), relative Serre vanishing gives the surjection onto
the first two jets along \(P\). Choose a rational function
\(t\) which is regular near \(P\), vanishes exactly to
order one there and has unit first derivative at \(x\).
Equivalently, in a neighbourhood of \(x\), the zero divisor
of \(t\) is exactly \(P\). Such a choice may be seen by
prescribing \(t\) modulo the ideal of \(2P\), using
\(H^1(X,\mathcal O_X(nD-2P))=0\). Its other zeros and
poles have reductions away from \(x\).

Normalize the actual proper curve \(X\) in
\[
                    K(X)(z),\qquad z^e=t.                 \tag{SB.1}
\]
The valuation at \(P_K\) makes this extension connected
of degree \(e\). After adjoining \(\mu_e\), it is cyclic.
At the unique closed point \(x'\) over \(x\), the
normalization is already regular and smooth over \(R\):
\[
 \widehat{\mathcal O}_{X,x}=R[[t]],\qquad
 R[[t]][z]/(z^e-t)=R[[z]].                               \tag{SB.2}
\]
Here the second expression means its completed local ring.
The finite algebra before completion is regular at that
point as well. The special inverse image of \(x\) has
the single underlying point \(z=0\).

After a further finite extension choose a semistable model
\(X'\) of the new generic curve which retains this smooth
neighbourhood. We spell out the preservation. The generic
residue extension on the original special component is
\(z^{p^a}=\overline t\). The differential of
\(\overline t\) is nonzero at \(x\), so this extension
is purely inseparable of the full degree \(p^a\). Its
smooth normalization has the same positive genus as \(X_s\)
and survives stabilization. In a common regular model,
all exceptional components over the already regular smooth
point \(x'\) arise by point blowups and are rational.
They form a tree attached to that retained component at
one point. There are no branch or auxiliary marked sections
there. Every such rational tree contracts in stabilization,
returning precisely the smooth neighbourhood (SB.2).
One may add auxiliary sections away from \(x'\) to use
pointed stable models. Resolution and contraction here use
the ordinary regular-model comparisons of
\cite[\S9.2 and \S10.3]{fghLiu}; no assertion that a
finite morphism extends finitely between arbitrary stable
models is being made.

\paragraph{The horizontal branch is killed only on the generic fibre.}
Take any connected component of the normalization of the
generic pullback to \(X'_K\), and denote it by \(Y'_K\).
Its cover of \(X'_K\) is Galois with a subgroup of \(G\),
hence with a \(p\)-group. At the horizontal valuation
\(P_K\), (SB.1) has ramification index \(e\).
The characteristic-zero discrete-valuation Abhyankar lemma
therefore kills the original inertia after this base change;
see the Stacks Project, Tag~0BRM. Thus \(Y'_K\to X'_K\)
is finite \'etale throughout the generic neighbourhood
above \(x'\). All its remaining branch sections have
reductions away from \(x'\), since their images in \(X_s\)
are different from \(x\).

This is an application at the horizontal valuation, whose
residue field has characteristic zero. It does \emph{not}
claim that the mixed-characteristic surface normalization
is \'etale or regular after base change. Vertical ramification
and singularities in that normalization may remain.
They are exactly what Raynaud's nonmarked smooth-point
theorem permits.

Exclude \(x'\) from the new branch marking and extend the
actual generic cover to a pointed semistable covering of
\(X'\), using Yang Proposition~1.6. Yang Theorem~1.9
then gives
\[
                    f(Y'_{x'})=0.                        \tag{SB.3}
\]
In particular every positive-genus component in this
vertical fibre has \(p\)-rank zero. Adding or removing
auxiliary marks changes only rational exceptional components.

\paragraph{Actual domination of each old vertical component.}
Let \(B\) be a positive-genus component of the original
semistable source which maps to \(x\). The chosen connected
generic pullback has a finite surjective map to the original
generic source. Resolve the rational map between their proper
models after resolving the nodal surfaces. This requires
only point blowups of a regular source surface; all newly
exceptional curves are rational. The resulting proper
generically finite morphism is surjective, so some component
\(B'\) dominates \(B\). A rational exceptional curve
cannot dominate a positive-genus smooth curve, even
inseparably, by L\"uroth's theorem. Hence \(B'\) is the
strict transform of a positive-genus component already
present in the semistable model of \(Y'\).

The diagram of actual function fields commutes with (SB.1).
Since \(B\) maps to \(x\), the image of \(B'\) in
\(X'\) is the unique point \(x'\) above \(x\).
Thus (SB.3) applies to \(B'\). The induced finite morphism
between the smooth normalizations of \(B'\) and \(B\)
forces \(f(B)\le f(B')=0\). For its inseparable part,
\(p\)-rank is invariant; for its separable part, the
Jacobian norm is surjective and \(p\)-rank cannot increase
on an abelian quotient. This proves the local assertion.
\end{proof}

\paragraph{Candidate global consequence at a rank-zero base.}
Under the same mixed-characteristic hypotheses, suppose
\(f(X_s)=0\), and the actual \(p\)-group cover is
branched at most at one \(R\)-section \(P\).
Then its Jacobian has potential good reduction of
special \(p\)-rank zero.

\begin{proof}[Candidate proof of the global consequence]
Potential good Jacobian reduction, equivalently a tree
for the semistable dual graph, follows from the ramified
nilpotent-cover theorem of Lehr--Matignon,
Theorem~6.13 \cite{fghLehrMatignon2005}. Its hypotheses
are the good base and the single \'etale horizontal branch
section. The local assertion just proved treats every
component contracted to \(P_s\); Raynaud's nonmarked
theorem treats components contracted to other points.

It remains to treat a component dominating \(X_s\).
Normalize \(X\) in the actual covering function field
and restrict to the regular open surface \(X\setminus P\).
Purity places its ramification support, if nonempty, on
the single smooth irreducible vertical divisor
\(X_s\setminus\{P_s\}\). Apply Raynaud,
Proposition~1(ii), \emph{on this open surface only}
\cite{fghRaynaud1990Pgroups}. The horizontal and vertical branch
divisors meet at \(P_s\), so the proposition is not
being applied across that crossing.

For each main component, its inertia is constant on this
open part. Quotient by it gives a connected finite \'etale
cover of \(X_s\setminus\{P_s\}\) with group \(D/I\),
where \(D\) is its decomposition group. The original
residue component maps radicially to that quotient.
After smooth proper normalization, the latter is an actual
connected \(p\)-group cover \(B_0\to X_s\) branched
at most at \(P_s\). This also holds if the vertical
ramification support is empty, when the cover is already
\'{e}tale there. The finite radicial map preserves
\(p\)-rank; it need not have degree exactly \(|I|\).

Write \(d=|D/I|\), and let \(r\ge1\) be the number
of points above \(P_s\). Deuring--Shafarevich gives
\[
 f(B_0)-1=d(0-1)+\sum_{Q\mid P_s}(e_Q-1)
              =-d+(d-r)=-r.                            \tag{SB.4}
\]
The same formula includes the unramified case, with
\(r=d\). Since \(f(B_0)\ge0\), necessarily \(r=1\)
and \(f(B_0)=0\). Every main component therefore has
rank zero. All component ranks and the graph Betti number
are zero, proving the stated Jacobian conclusion.
\end{proof}

\paragraph{Proposed arithmetic application, pending the preceding audits.}
For the fixed elliptic curve \(E/\mathbb Q(i)\), take
\(p=2\) or any odd \(p\equiv3\pmod4\). The curve
has potential good supersingular reduction at \(p\).
If an actual finite geometric \(p\)-group closure \(G\)
has cyclic derived subgroup, then both assertions of
Task~A.5 follow at every coefficient prime, including all
actual intermediates. The same is true after an additional
abelian factor of order prime to \(p\), via the actual
connected isogeny pullback and rational-path preservation
already proved in A.66/A.76.

Indeed the two actual handle generators generate \(G\).
Their commutator normally generates \(G'\), and every
subgroup of the cyclic group \(G'\) is characteristic
in \(G'\). Thus the actual cusp inertia is exactly
\(I=G'\). The proper quotient \(C/I\) is an actual
elliptic isogeny cover of \(E\), with kernel \(G/I\).
If that kernel has exponent \(d\), pulling back its
torsion divisors through \(C\to C/I\) kills every
cusp difference by \(d|I|\), a power of \(p\).
The proposed global consequence gives potential good
Jacobian reduction of rank zero. For every proper
subgroup \(H<G\), normality of \(I\) gives cusp count
\([G:HI]\ge p\), since \(I=G'\subset\Phi(G)\).
Actual norm maps preserve the rank-zero conclusion and
the cusp-divisor relation on every intermediate. A.41
and A.42 give the rational companion paths and all-prime
fixed-inertia assertions, and A.59 gives all normalizers.
The abelian and original genus-one cases are A.40.
This does not settle arbitrary \(p\)-group fixedness:
without normal cyclic inertia, the required cusp-divisor
or rational-path comparison remains a separate question.
